%!TEX root = paper.tex

\section{Introduction}

Module extraction is the task of computing, given an ontology
$\T$ and a signature of interest $\Sigma$, a (preferably small) subset $\M$ of
$\T$ (a module) that preserves all relevant entailments
in $\T$ over the set of symbols $\Sigma$. 
Such an $\M$ is indistinguishable from $\T$ w.r.t.\ $\Sigma$, 
and $\T$ can be safely replaced with $\M$
in applications of $\T$ that use only the symbols in $\Sigma$. 
%The utility of such a replacement is usually dependent on the size 
%of $\M$; ideally, it should be of minimal size.

Module extraction has received a great deal of attention in recent years
\cite{DBLP:series/lncs/5445,GrauHKS08JAIR,DBLP:conf/www/SeidenbergR06,DBLP:journals/ai/KontchakovWZ10,DBLP:conf/dlog/GatensKW14,DBLP:conf/ijcai/VescovoPSS11,DBLP:conf/lpar/NortjeBM13}, 
and modules have found a wide range of applications, including
ontology reuse \cite{GrauHKS08JAIR,DBLP:conf/esws/Jimenez-RuizGSSL08},  matching \cite{DBLP:conf/semweb/Jimenez-RuizG11},
debugging \cite{DBLP:conf/aswc/SuntisrivarapornQJH08,LudwigORE} and
classification \cite{DBLP:conf/semweb/RomeroGH12,DBLP:conf/ore/TsarkovP12,DBLP:journals/jar/GrauHKS10}.
   
The preservation of relevant entailments is formalised via
\emph{inseparability relations}.
The strongest notion is \emph{model} inseparability, which requires
that it must be possible
to turn any
model of $\M$ into a model of $\T$
by (re-)interpreting only the symbols outside $\Sigma$; such an 
$\M$ preserves all second-order
entailments of $\T$ w.r.t.\ $\Sigma$ \cite{DBLP:journals/ai/KonevL0W13}.
A weaker and more flexible notion is \emph{deductive} inseparability,
which requires only that
$\T$ and $\M$ 
entail the same $\Sigma$-formulas \emph{in a given 
query language}.
Unfortunately, the decision problems
associated with module extraction are invariably of high complexity, 
and often undecidable.
For model inseparability, checking whether $\M$ is a $\Sigma$-module
in $\T$ 
is undecidable  even if $\T$ is restricted to be
in the description logic (DL) $\mathcal{EL}$,
for which standard reasoning is tractable.
For deductive inseparability, the problem is typically decidable for lightweight DLs and ``reasonable'' query languages, albeit of high worst-case complexity; e.g., the problem is already
\textsc{ExpTime}-hard for $\mathcal{EL}$ if we consider concept inclusions as
the query language \cite{DBLP:journals/jsc/LutzW10}.
Practical algorithms that ensure minimality of
the extracted modules are known only for 
acyclic $\mathcal{ELI}$ \cite{DBLP:journals/ai/KonevL0W13}
and DL-Lite \cite{DBLP:journals/ai/KontchakovWZ10}.

Practical module extraction techniques are typically based on
sound approximations: they ensure that the extracted fragment
$\M$ is a module (i.e., inseparable from $\T$ w.r.t.\ $\Sigma$),
but they give no minimality guarantee. The most popular
such techniques are based on a family of polynomially checkable
conditions called syntactic 
locality \cite{DBLP:conf/www/GrauHKS07,GrauHKS08JAIR,WhichKindOfModule}; in
particular, $\bot$-locality  and 
$\top\!\bot\!^*$-locality.  
Each locality-based module $\M$ enjoys a number of desirable properties
for applications:
\begin{inparaenum}[\it (i)]
 \item it is model inseparable from $\T$;
 \item it is \emph{depleting}, in the sense that $\T \setminus \M$
 is inseparable from the empty ontology w.r.t.\ $\Sigma$; 
 \item it contains all justifications  (a.k.a.\ explanations) in $\T$
 of every $\Sigma$-formula entailed by $\T$; and
\item last but not least, it can be computed
efficiently, even for very expressive
ontology languages.
\end{inparaenum}

Locality-based techniques are easy to implement, and
surprisingly effective in practice. Their main drawback is that the
extracted modules can be rather large, which limits their 
usefulness in some applications \cite{VescovoKPS0T13}.
One way to address this issue is to develop techniques that more closely approximate minimal modules while still preserving properties \emph{(i)}--\emph{(iii)}.
Efforts in this direction have confirmed that locality-based modules can
be far from optimal in practice \cite{DBLP:conf/dlog/GatensKW14};
however, these techniques
apply only to rather restricted ontology languages and utilise algorithms with high worst-case
\mbox{complexity}.

Another approach to computing smaller modules is to weaken 
properties \emph{(i)}--\emph{(iii)}, which are stronger
than is required in many applications.
In particular, model inseparability (property \emph{(i)}) is
a very strong condition, and
deductive inseparability would usually suffice,
with the query language determining which kinds of 
consequence are preserved; in modular classification, for example,
only atomic concept inclusions need to be preserved.
However, 
%with the single exception of \cite{DBLP:conf/lpar/NortjeBM13},
all practical module extraction techniques that are applicable
to expressive ontology languages yield modules satisfying
all three properties, and  hence potentially 
much larger than they need to be.

%% Present our solution. 
 
In this paper,   
we propose
a technique that reduces module 
extraction to a reasoning problem in datalog.
The connection between module extraction and datalog 
was first observed in \cite{DBLP:conf/esws/Suntisrivaraporn08}, where it was shown that 
locality $\bot$-module extraction for $\mathcal{EL}$ ontologies 
could be reduced to propositional datalog
reasoning.
Our approach takes this connection much farther by
generalising both locality-based and
reachability-based \cite{DBLP:conf/lpar/NortjeBM13} modules
for expressive ontology languages in an elegant way.
A key 
distinguishing feature of our technique is that it
can extract deductively inseparable modules, with the
query language tailored to the requirements of the
application at hand, 
which allows
us to relax Property \emph{(i)} and extract
significantly smaller modules. In all cases our modules preserve
the nice features of locality: they are
widely applicable (even beyond DLs), they
can be efficiently computed,  
they
are depleting  (Property \emph{(ii)}) and they
 preserve all justifications of relevant entailments
(Property \emph{(iii)}). 

We have implemented our approach
using the RDFox datalog engine \cite{DBLP:conf/aaai/MotikNPHO14}. 
Our proof of concept
evaluation shows that
module size consistently decreases as we consider weaker inseparability
relations, which could significantly improve the usefulness of modules in applications.

All our proofs are deferred to the appendix.
%  an extended version of the
% paper available online at 
% \texttt{arXiv:1411.5313}.
